\documentclass{article}
\usepackage{amsmath}
\usepackage{microtype}

\begin{document}

The vowel that shifted most dramatically in connected speech was /\ae/, showing a marked centralizing drift in \(F_1\) and \(F_2\) compared to its isolated /hVd/ baseline. In continuous speech, shorter duration prevented the tongue and jaw from reaching extreme peripheral positions, leading to articulatory target undershoot. Consequently, \(F_1\) decreased due to reduced jaw opening, while \(F_2\) shifted toward a mid-range frequency typical of schwa-like central reduction. In contrast, vowels in isolated /hVd/ contexts retained higher acoustic clarity and distinct formant targets due to hyperarticulated citation forms.

\end{document}